% SOURCE_AUTHORITY: sga5_fr_workpass.tex, lines 13818--13845 (LF-normalized unit).
% SOURCE_SHA256: 791F4EFFC5E02832D5D77ED03518C8156D6F07E4C8238B03545DB93D883FBB28
\section*{6. Aplicação a uma fórmula de Lefschetz}

\subsection*{6.1}
Conservaremos as notações e condições do número 1\footnote{$X^f$ denota o conjunto dos pontos fixos fechados de $X$.}. Sejam, além disso, $A$ uma $\La_n$-álgebra comutativa noetheriana, $F$ um feixe (para a topologia étale) construtível de $A$-módulos perfeitos sobre $X$, e
\[
u : f^*(F) \longrightarrow F
\]
um homomorfismo de feixes de $A$-módulos.

Em particular, obtém-se um endomorfismo do par $(X, F)$:
\[
(u, f) : (X, F) \longrightarrow (X, F) \quad ,
\]
e, consequentemente, um endomorfismo de complexos de $A$-módulos:
\[
\R\Gamma_X(u, f) : \R\Gamma_X(F) \longrightarrow \R\Gamma_X(F) \quad .
\]

Como, por hipótese, o esquema $X$ é próprio sobre $k$ e o feixe $F$ é construtível e de torção, o teorema de finitude (S.G.A. 4. XVII 5) assegura que o complexo $\R\Gamma_X(F)$ é um complexo perfeito de $A$-módulos.

O traço $\Tr_A(\R\Gamma_X(u, f))$ está, portanto, definido:
\[
\Tr_A(\R\Gamma_X(u, f)) \in A \quad .
\]

Nosso objetivo será dar, sob algumas hipóteses suplementares, uma fórmula explícita para $\Tr_A(\R\Gamma_X(u, f))$ em termos de invariantes locais associados aos pontos fixos de $f$.
